% Part: computability
% Chapter: computability-theory
% Section: computable-sets

\documentclass[../../../include/open-logic-section]{subfiles}

\begin{document}

\olfileid[pt]{cmp}{thy}{cps}
\olsection{Conjuntos computáveis}

Podemos estender a noção de computabilidade das funções computáveis
aos conjuntos computáveis:

\begin{defn}
  Seja $S$ um conjunto de números naturais. Então $S$ é
  \emph{computável} se, e somente se, sua função
  característica~$\Char{S}$ o é. Em outras palavras, $S$~é computável
  se, e somente se, a função
\[
\Char{S}(x) =
\begin{cases}
1 & \text{se $x \in S$} \\
0 & \text{caso contrário}
\end{cases}
\]
é computável. De modo semelhante, uma relação
$R(x_0, \dots, x_{k-1})$ é computável se, e somente se, sua função
característica é computável.

Conjuntos e relações computáveis também são chamados de
\emph{decidíveis}.
\end{defn}

\begin{explain}
Observe que agora temos várias noções de computabilidade: para funções
parciais, para funções e para conjuntos. Não as confunda!{}
\iftag{TMs}{A máquina de Turing que calcula uma função parcial devolve
  a saída da função para os valores de entrada nos quais a função está
  definida; a máquina de Turing que calcula um conjunto devolve $1$
  ou~$0$, depois de decidir se o valor de entrada pertence ou não ao
  conjunto.}{}
\end{explain}

\end{document}

